\documentclass[../../main.tex]{subfiles}
\begin{document}
\section{Family 1: classical needle localization}
\label{sec:family-needles}

\paragraph{Object followed.}
One-dimensional restrictions of $\mu$ to segments --- \emph{needles} --- carrying a log-concave
weight.

\paragraph{What it buys.}
The KLS localization lemma reduces certain $n$-dimensional integral inequalities to inequalities
on a segment \cite{KannanLovaszSimonovits1995}. Schematically: given two integral constraints
\[
  \int_{\R^n}g_1\dd x\ge0,\qquad \int_{\R^n}g_2\dd x\ge0,
\]
it produces points $a,b\in\R^n$ and an affine $\ell>0$ such that both constraints remain valid for
the one-dimensional measure
\begin{equation}
  \ell(t)^{n-1}\dd t,\qquad 0\le t\le1,
  \label{eq:needle-measure}
\end{equation}
carried on the segment $x(t)=(1-t)a+tb$. The weight $\ell^{n-1}$ is log-concave, so the reduced
problem is a one-dimensional log-concave problem, and one-dimensional log-concave isoperimetry is
completely understood: for a log-concave $\nu$ on $\R$ with coordinate $T$,
\begin{equation}
  h_\nu\asymp\frac1{\sqrt{\Var_\nu T}} .
  \label{eq:one-dim-isoperimetry}
\end{equation}

\paragraph{Sharpest result.}
Combining \eqref{eq:one-dim-isoperimetry} with localization gives
\begin{equation}
  h_\mu\gtrsim\frac1{\sqrt{\Tr\Cov\mu}},
  \label{eq:needle-bound}
\end{equation}
which in isotropic position ($\Tr\Cov\mu=n$) is the classical bound
\[
  \PsiKLS_n\lesssim\sqrt n ,
\]
the 1995 entry of the history table in Section~\ref{subsec:kls-status}.

\paragraph{The precise missing estimate.}
A covariance-sensitive decomposition: one whose needles inherit the \emph{operator} covariance of
$\mu$, not merely one or two scalar integral constraints.

\paragraph{Why it stalls.}
The obstruction is exact, and it is a counting statement. The localization construction preserves
one or two scalar integral constraints. Isotropy is $n(n+1)/2$ constraints. Nothing forces an
individual needle to be even approximately isotropic, and in fact a single needle can have
variance of order $n$ while the original measure has $\Cov\mu=I$. This is precisely why
\eqref{eq:needle-bound} sees $\Tr\Cov\mu$ --- a sum of $n$ eigenvalues --- where the affine
conjecture \eqref{eq:kls-affine} asks for $\norm{\Cov\mu}_\op$, the largest one alone. The gap
between $\Tr$ and $\norm{\cdot}_\op$ in \eqref{eq:needle-bound} \emph{is} the factor $\sqrt n$.

Klartag's needle decomposition, built from transport rays rather than from the bisection
construction, is substantially more geometric and can be arranged so that a chosen mean-zero
function remains mean-zero on almost every needle \cite{Klartag2014Needle}. That is a genuine
strengthening --- it is one preserved constraint chosen adaptively rather than arbitrarily --- but
it still does not deliver uniformly bounded conditional covariance on the needles.

\begin{remark}[What a deterministic route would have to supply]
\label{rem:needle-requirement}
A successful deterministic localization proof of KLS would need a genuinely new
covariance-sensitive decomposition, not merely sharper one-dimensional inequalities. The
one-dimensional theory in \eqref{eq:one-dim-isoperimetry} is already sharp; there is no slack left
to extract there. This is the structural reason the field moved to the stochastic construction of
Section~\ref{sec:family-sl}, which preserves covariance information by construction --- at the
price of preserving it only in law, along a random path.
\end{remark}

\paragraph{How this family feeds the live routes.}
It does not, directly, and that is worth saying plainly: no route in Parts~III or~IV decomposes
$\mu$ into one-dimensional pieces. Route~C (Section~\ref{sec:moment-map-cmh}) is the only
deterministic route here, but it works in the moment-map coordinates of
Section~\ref{sec:family-moment-map} rather than by needle decomposition. What the family
contributes is the diagnosis --- global isotropy is not inherited needle by needle --- which is
the first instance of the fixed-versus-adaptive difficulty of Section~\ref{sec:kls-remaining}.
\end{document}
